Event data on political violence and protest are now available at fine resolution
across much of the world. The Armed Conflict Location \& Event Data Project
\citep[ACLED;][]{raleigh2010acled} records the date, place and type of individual
battles, attacks on civilians, riots and demonstrations, and such data underpin
early-warning systems for armed conflict \citep{hegre2019views, hegre2025views} and a
growing statistical literature on where and when violence occurs. Recorded events are
rare and strongly clustered. Most cells of a $0.5^\circ$ grid record no event in a
given week, a few record events in most weeks, and the strongest single predictor of
an event is recent events in the same place \citep{weidmann2010predicting,
mueller2022hard, vestby2022predicting}. Spatial statistical models represent this
clustering with latent fields \citep{zammitmangion2012afghan, python2019bayesian,
python2021predicting, egbon2024spatio}, and point-process models with self-excitation
\citep{mohler2011self, clark2018selfexciting, reinhart2018review}. Both have mostly
modelled one type of event at a time.

Different types of events are, however, closely linked. Protests and riots are often
followed by armed conflict in the same area, and an influential reading of this
association is escalation. Repression of protest can radicalise movements and push
them towards violence \citep{ritter2016preventing, ryckman2020turn, rod2023protest},
and the outcome of nonviolent campaigns shapes whether violent ones follow
\citep{chenoweth2017violent}. Forecasting models therefore use the recent history of
unrest, and its spatial lags, as predictors of conflict \citep{rod2025predicting}, and
diffusion models add spatial lags of past events to the predictor
\citep{buhaug2008contagion, racek2025diffusion}. A second reading is that both kinds of
events respond to the same latent episodes: a period of political crisis in a region,
a contested election, a collapse of local authority, or simply a period in which media
and local partners pay attention to a place, raising the recorded counts of every
type of event \citep{weidmann2016closer, hellmeier2018spotlight, croicu2022reporting}.
Under this reading, recent unrest predicts conflict because it reveals the current
state of the episode, not because it causes violence.

The two readings are hard to separate with the models that are usually fitted.
A regression of conflict on lagged unrest, with or without a latent field for
conflict, attributes to the lag coefficient whatever part of the shared episode the
conflict data alone cannot reveal. This is a familiar problem in another guise: in
panel models, state dependence is confounded with persistent heterogeneity unless the
latter is modelled \citep{heckman1981heterogeneity, wooldridge2005initial}, and a
latent effect estimated from rare binary outcomes is too noisy to absorb it fully. The
remedy we pursue is to model the two outcomes jointly, so that unrest data inform the
latent episode directly rather than only through a covariate.

We propose a joint latent Gaussian model \citep{rue2009inla} for the weekly
occurrence of armed conflict and of social unrest in cells of a $0.5^\circ$ grid. Each
outcome has its own fixed effects, cross-lagged effects of both outcomes' recent
histories in the cell and its neighbours, country and cell effects and a static
Mat\'ern field. The two outcomes share a space-time field, an AR(1) process of
Mat\'ern innovations \citep{cameletti2013pm10} that enters conflict with loading one
and unrest with an estimated loading, in the manner of a linear model of
coregionalisation \citep{gelfand2004nonstationary}. The coefficient of recent unrest
on conflict is then an association conditional on shared episodes, which we call
escalation conditional on shared episodes; an extended model adds a space-time field
for unrest alone, so that episodes of unrest without conflict are not forced into the
shared field. Fields are defined through the SPDE representation of Mat\'ern fields
\citep{lindgren2011spde}, discretised by finite volumes directly on the
longitude--latitude grid on the sphere, with sea barriers \citep{bakka2019barrier} so
that correlation does not leak across narrow seas. Inference is an empirical Bayes
Laplace approximation in the spirit of INLA and TMB \citep{rue2009inla,
kristensen2016tmb}. Because the regions analysed contain millions of cell-weeks per
outcome, we estimate the hyperparameters on a case-control subsample of the non-events
with an exact offset \citep{prentice1979logistic}, and compute the latent posterior on
all data; this reduces the cost of a fit several-fold without changing the estimates
or the predictions materially.

The paper makes three contributions. First, a simulation study shows that univariate
models, including models with their own space-time field, report escalation that does
not exist when the two outcomes share latent episodes, while the joint model recovers
the escalation coefficient with close to nominal coverage, also when unrest has
episodes of its own. Second, in ACLED data for West Africa and the Sahel and for
Africa and the Middle East, 2018--2024, the coefficient of recent unrest on conflict
falls from 0.13--0.15 in univariate models to about 0.10 once both shared and
unrest-specific episodes are modelled, and to about 0.05 with a single shared field:
shared episodes account for between a quarter and two thirds of the univariate
association. Escalation at the short, local scale that the model can resolve is real
but smaller than univariate models suggest. Third, we examine what the joint model
adds to prediction. When both outcomes are observed, it cannot forecast either one
better than univariate models that use both histories as covariates, and it does not.
When one outcome is withheld in an area while the other is observed, as when one type
of event goes unreported, the joint model predicts the withheld outcome better, in
simulations and, for conflict, in the data; the gain comes mostly from the weeks in
which the other outcome is observed and fades in forecasts.

Section~\ref{sec:data} describes the data and Section~\ref{sec:methods} the model and
its inference. Section~\ref{sec:simulation} reports the simulation studies and
Section~\ref{sec:application} the application. Section~\ref{sec:discussion}
discusses the interpretation of the shared episodes, the limits of the analysis and
extensions. The Supplementary Material contains the construction of the data,
computational details, further simulation studies of the univariate building blocks,
and further results.
